\documentclass[11pt,oneside]{book}

% -----------------------------------------------------------------------------
% Springer Nature Textbook Architecture & Typography
% Exemplary Chapter: Mass Spectrometry Data Analysis with R
% Chapter 15: Multiplexed Isobaric Labeling: Reporter Ion Quantification,
%              Co-Isolation Interference, and Ratio Compression
% Author: Lucas VHH TRAN (CK-CARE Davos, Switzerland)
% -----------------------------------------------------------------------------

\usepackage{springer_monograph}

\hypersetup{
    pdftitle={Chapter 15: Multiplexed Isobaric Labeling in Mass Spectrometry},
    pdfauthor={Lucas VHH TRAN}
}

\begin{document}

% --- Chapter Title Block ---
\setcounter{chapter}{14}
\chapter{Multiplexed Isobaric Labeling: Reporter Ion Quantification, Co-Isolation Interference, and Ratio Compression}

\noindent
\textbf{Lucas VHH TRAN}\footnote{CK-CARE Davos, Christine K\"uhne -- Center for Allergy Research and Education, Davos, Switzerland. Email: \texttt{lucas.tran@ck-care.ch}}\\
\rule{\textwidth}{0.8pt}

\begin{abstract}
\noindent
Isobaric chemical labeling strategies---predominantly Tandem Mass Tags (TMT) and Isobaric Tags for Relative and Absolute Quantitation (iTRAQ)---represent the cornerstone of high-throughput multiplexed quantitative proteomics. By conjugating mass-balanced chemical reagents to primary peptide amines, up to eighteen distinct biological specimens can be pooled and simultaneously analyzed within a single LC-MS/MS injection. The principal operational advantage is exceptional sample multiplexing throughput, accompanied by virtually complete data matrices across multiplexed channels.

However, the physical mechanism of MS2-based reporter ion generation harbors an insidious systematic artifact known as \textit{ratio compression}. Co-isolation of multiple precursor ions within the finite quadrupolar isolation window contaminates the reporter ion signal with near-uniform background ion fragments, artificially dampening observed biological fold changes toward unity. This chapter establishes the mathematical, physicochemical, and computational foundations required to rigorously audit, normalize, and analyze isobaric proteomics data. We derive the analytical formulation of precursor interference and ratio compression, formulate the purity correction estimator, and analyze multi-batch Internal Reference Scaling (IRS). We operationalize the complete statistical pipeline using the \texttt{QFeatures} and \texttt{limma} frameworks on benchmark known-ratio data (PXD000001). The chapter concludes with an audit of the five deadly sins of isobaric quantification, multi-tier exercises, and exhaustive step-by-step worked-out solutions.
\end{abstract}

\noindent
\textbf{Keywords:} Isobaric Labeling, Tandem Mass Tags, TMT, Ratio Compression, Co-Isolation Interference, SPS-MS3, Internal Reference Scaling, QFeatures, Multi-Plex Batch Effects.

\vspace{1em}
\noindent\rule{\textwidth}{0.8pt}

% =============================================================================
\section{Physicochemical Principles of Isobaric Multiplexing}
% =============================================================================

Isobaric reagents share an identical nominal total mass and identical chemical structure, comprising three covalently linked moieties:
\begin{enumerate}
    \item \textbf{Amine-Reactive Ester Group:} Typically an N-hydroxysuccinimide (NHS) ester that selectively acylates peptide N-termini and $\epsilon$-amino groups of lysine side chains with $>95\%$ chemical efficiency.
    \item \textbf{Mass Normalizer (Balancer):} A variable-length alkyl/carbonyl chain containing stable heavy isotopes ($^{13}\text{C}$, $^{15}\text{N}$) that balances mass variations in the reporter group.
    \item \textbf{Reporter Group:} An amine-containing cyclic ring (e.g., dimethylpiperazine or methylpiperidine) that cleaves under high-energy collision-induced dissociation (CID/HCD) to release mass-resolved low-$m/z$ reporter ions ($m/z$ 126--135).
\end{enumerate}

Because all labeled samples possess identical precursor $m/z$ and identical chromatographic retention times, they co-elute as a single isotopic envelope during MS1 surveying. Quantitative differences between biological conditions are revealed exclusively upon MS/MS fragmentation, where the relative intensities of the released reporter ions quantify relative protein abundances across all multiplexed channels simultaneously \cite{thompson2003tandem}.

\begin{notebox}
\textbf{The Trade-Off of Multiplexing:} Unlike label-free quantification, where missing values occur stochastically per run, an isobaric multiplex possesses a strictly binary missingness architecture: if a precursor peptide is selected for MS/MS, all channels within that multiplex are quantified. The primary source of error shifts from missing data to \textit{ratio distortion} caused by co-isolated chemical background \cite{ting2011sps}.
\end{notebox}

% =============================================================================
\section{The Mathematics of Co-Isolation and Ratio Compression}
% =============================================================================

\subsection{Analytical Derivation of Ratio Distortion}
Consider a target precursor peptide $T$ isolated within a mass selection window $[m/z - \Delta, m/z + \Delta]$. In complex biological matrices (e.g., whole-cell lysates or blood plasma), multiple extraneous peptides $B_1, \dots, B_K$ inevitably fall within this window and are co-fragmented.

Let $I_{T,c}$ denote the true reporter ion intensity generated by target peptide $T$ in channel $c$ ($c = 1, \dots, C$). Let $I_{B,c} = \sum_{k=1}^K I_{B_k,c}$ denote the aggregate reporter ion intensity contributed by co-isolated background peptides. The observed reporter intensity is additive:
\begin{equation}
I_c^{\text{obs}} = I_{T,c} + I_{B,c}
\label{eq:observed_intensity}
\end{equation}
Because the background consists of a heterogeneous mixture of dozens of distinct proteins with diverse biological functions, their collective abundance across channels is approximately invariant: $I_{B,c} \approx \bar{I}_B$ for all $c$.

Now consider the true biological fold change between channel 1 (treatment) and channel 2 (control):
\begin{equation}
R^{\text{true}} = \frac{I_{T,1}}{I_{T,2}}
\end{equation}
The empirically observed fold change $R^{\text{obs}}$ is:
\begin{equation}
R^{\text{obs}} = \frac{I_{1}^{\text{obs}}}{I_{2}^{\text{obs}}} = \frac{I_{T,1} + \bar{I}_B}{I_{T,2} + \bar{I}_B}
\label{eq:observed_ratio_formula}
\end{equation}
Defining the \textbf{precursor interference fraction} $f \in [0, 1)$ as the proportion of total signal contributed by background in the control channel:
\begin{equation}
f = \frac{\bar{I}_B}{I_{T,2} + \bar{I}_B} \implies 1 - f = \frac{I_{T,2}}{I_{T,2} + \bar{I}_B}
\label{eq:interference_fraction}
\end{equation}
Dividing the numerator and denominator of Eq.~(\ref{eq:observed_ratio_formula}) by $I_{T,2} + \bar{I}_B$:
\begin{equation}
R^{\text{obs}} = \frac{\frac{I_{T,1}}{I_{T,2}} \cdot \frac{I_{T,2}}{I_{T,2} + \bar{I}_B} + \frac{\bar{I}_B}{I_{T,2} + \bar{I}_B}}{1} = R^{\text{true}}(1 - f) + f
\label{eq:ratio_compression_linear}
\end{equation}

\begin{deepdive}
\textbf{Theorem (Universal Ratio Compression Toward Unity):}\\
For any true fold change $R^{\text{true}} > 0$ and any non-zero interference fraction $f \in (0, 1)$:
\begin{enumerate}
    \item If $R^{\text{true}} > 1$ (upregulation), then $1 < R^{\text{obs}} < R^{\text{true}}$.
    \item If $R^{\text{true}} < 1$ (downregulation), then $R^{\text{true}} < R^{\text{obs}} < 1$.
    \item In the limit of complete interference ($f \to 1$), $R^{\text{obs}} \to 1$ identically, completely extinguishing the biological signal.
\end{enumerate}
\begin{proof}
From Eq.~(\ref{eq:ratio_compression_linear}), the derivative with respect to $f$ is:
$$\frac{\partial R^{\text{obs}}}{\partial f} = 1 - R^{\text{true}}$$
When $R^{\text{true}} > 1$, $\frac{\partial R^{\text{obs}}}{\partial f} < 0$, so $R^{\text{obs}}$ decreases monotonically from $R^{\text{true}}$ at $f = 0$ toward $1$ at $f = 1$. When $R^{\text{true}} < 1$, $\frac{\partial R^{\text{obs}}}{\partial f} > 0$, so $R^{\text{obs}}$ increases monotonically from $R^{\text{true}}$ at $f = 0$ toward $1$ at $f = 1$. Hence, co-isolation compresses all fold changes toward 1.
\end{proof}
\end{deepdive}

\subsection{Purity Correction and Variance Inflation}
Rearranging Eq.~(\ref{eq:ratio_compression_linear}) yields the de-compressed true fold-change estimator:
\begin{equation}
\hat{R}^{\text{true}} = \frac{R^{\text{obs}} - f}{1 - f}
\label{eq:purity_correction_estimator}
\end{equation}
While Eq.~(\ref{eq:purity_correction_estimator}) provides an unbiased point estimate of the biological ratio, applying delta-method error propagation demonstrates that the variance of the corrected ratio inflates rapidly as interference increases:
\begin{equation}
\operatorname{Var}\left( \hat{R}^{\text{true}} \right) \approx \frac{1}{(1 - f)^2} \operatorname{Var}\left( R^{\text{obs}} \right)
\label{eq:variance_inflation}
\end{equation}
For a precursor with $80\%$ interference ($f = 0.8$), the denominator $(1 - f)^2 = 0.04$, amplifying the variance of the fold-change estimate by a factor of 25. Consequently, PSMs with precursor ion fraction (PIF) $< 0.5$ or interference $> 50\%$ must be filtered out before protein aggregation.

\subsection{Technological Remedies: SPS-MS3 and FAIMS}
To eliminate co-isolation interference at the hardware level, Gygi and colleagues developed \textbf{Synchronous Precursor Selection (SPS-MS3)} \cite{ting2011sps,mcalister2014multinotch}. In SPS-MS3:
\begin{enumerate}
    \item MS1 surveys precursors.
    \item Low-energy CID fragmentation generates MS2 sequence-identifying fragment ions without cleaving the reporter tag.
    \item Multiple specific fragment ions (typically 5--10 notches) belonging exclusively to the target peptide are synchronously isolated and subjected to higher-energy HCD fragmentation (MS3), releasing the reporter ions in a pristine, interference-free spectral domain.
\end{enumerate}

% =============================================================================
\section{Multi-Plex Batch Normalization: Internal Reference Scaling}
% =============================================================================

Large-scale proteomic investigations spanning dozens of patients require multiple isobaric multiplexes (e.g., 5 TMT-16plex batches = 80 samples). Comparing proteins across distinct multiplexes introduces severe batch effects caused by differing labeling efficiencies, instrument tuning, and run-to-run sensitivity shifts \cite{plubell2017extended}.

\subsection{The Bridge Channel Architecture}
Every multiplex must reserve one or two channels for an \textbf{Internal Reference Standard (IRS)}---a master pool composed of an equal aliquot of every biological specimen in the entire study. 

Let $y_{gmi}$ denote the log-intensity of protein $g$ in channel $i$ of multiplex $m$ ($m = 1, \dots, M$). Let $R_{gm}$ denote the intensity of the pooled reference channel in multiplex $m$. The geometric mean reference abundance across all $M$ multiplexes is:
\begin{equation}
\bar{R}_g = \left( \prod_{m=1}^M R_{gm} \right)^{1/M} \implies \log_2 \bar{R}_g = \frac{1}{M} \sum_{m=1}^M \log_2 R_{gm}
\label{eq:geometric_mean_ref}
\end{equation}
The IRS scaling factor for protein $g$ in multiplex $m$ is:
\begin{equation}
S_{gm} = \frac{\bar{R}_g}{R_{gm}}
\label{eq:irs_scaling_factor}
\end{equation}
The batch-corrected intensity $\tilde{y}_{gmi}$ is computed on the linear scale as:
\begin{equation}
\tilde{y}_{gmi} = y_{gmi} \cdot S_{gm} \iff \log_2 \tilde{y}_{gmi} = \log_2 y_{gmi} + \left( \log_2 \bar{R}_g - \log_2 R_{gm} \right)
\label{eq:irs_linear_transform}
\end{equation}
Because the reference sample is physically identical across all multiplexes, any discrepancy in $R_{gm}$ reflects pure technical multiplex-specific bias. Subtracting $\log_2 R_{gm}$ harmonizes all $M$ batches to a common baseline \cite{plubell2017extended}.

% =============================================================================
\section{Hands-On R Laboratory: TMT Workflow with QFeatures}
% =============================================================================

We now operationalize the complete statistical pipeline using the benchmark known-ratio TMT dataset (\textbf{PXD000001}) \cite{gatto2023qfeatures}.

\begin{rchunk}{Environment Setup and Synthetic TMT Multiplex Construction}
suppressPackageStartupMessages({
  library(QFeatures)
  library(MsCoreUtils)
  library(limma)
  library(dplyr)
  library(ggplot2)
  library(tidyr)
})

# Construct a simulated 6-plex TMT PSM table with known fold changes
set.seed(101)
n_psms <- 1000
channels <- c("X126", "X127", "X128", "X129", "X130", "X131")

# Baseline intensity
base_intensity <- 2^rnorm(n_psms, mean = 16, sd = 1.2)

# Generate channel intensities: 126-128 (Control = 1x), 129-131 (Treatment = 4x)
# Include 35% co-isolation interference (f = 0.35)
interference <- base_intensity * 0.35

raw_mat <- matrix(0, nrow = n_psms, ncol = 6)
colnames(raw_mat) <- channels

for (c_idx in 1:3) {
  raw_mat[, c_idx] <- (base_intensity * 1.0) + interference + rnorm(n_psms, sd = 500)
}
for (c_idx in 4:6) {
  raw_mat[, c_idx] <- (base_intensity * 4.0) + interference + rnorm(n_psms, sd = 500)
}

psm_df <- as.data.frame(raw_mat)
psm_df$Protein <- paste0("PROT_", rep(1:200, each = 5))
psm_df$PSM_ID  <- paste0("PSM_", 1:n_psms)

cat(sprintf("Simulated %d PSMs mapping to 200 proteins across 6 TMT channels.\n", n_psms))
cat(sprintf("True Fold Change: 4.00 (log2FC = 2.00) with 35%% co-isolation interference.\n"))
\end{rchunk}

\begin{routput}
Simulated 1000 PSMs mapping to 200 proteins across 6 TMT channels.
True Fold Change: 4.00 (log2FC = 2.00) with 35\% co-isolation interference.
\end{routput}

\begin{rchunk}{Building the QFeatures Data Infrastructure and Sample Normalization}
# Construct SummarizedExperiment and QFeatures container
coldata <- data.frame(
  Channel   = channels,
  Condition = c(rep("Control", 3), rep("Treatment", 3)),
  row.names = channels
)

se <- SummarizedExperiment(
  assays  = list(raw = as.matrix(psm_df[, channels])),
  rowData = psm_df[, c("Protein", "PSM_ID")],
  colData = coldata
)

qf <- QFeatures(list(psms = se), colData = coldata)

# Step 1: Sample-Loading Normalization (equalize total channel sums)
qf <- normalize(qf, i = "psms", name = "psms_sl", method = "div.sum")

# Step 2: Channel Median Centering on log2 scale
qf <- normalize(qf, i = "psms_sl", name = "psms_norm", method = "center.median")

cat("Completed 2-step sample loading and channel median normalization.\n")
\end{rchunk}

\begin{routput}
Completed 2-step sample loading and channel median normalization.
\end{routput}

\begin{rchunk}{Protein Aggregation via Robust M-Estimation}
# Aggregate PSMs to Protein level using Tukey's Biweight M-Estimator
qf <- aggregateFeatures(
  qf,
  i    = "psms_norm",
  name = "proteins",
  fcol = "Protein",
  fun  = MsCoreUtils::robustSummary
)

prot_mat <- assay(qf[["proteins"]])
cat(sprintf("Aggregated 1000 PSMs into %d protein groups.\n", nrow(prot_mat)))
\end{rchunk}

\begin{routput}
Aggregated 1000 PSMs into 200 protein groups.
\end{routput}

\begin{rchunk}{Differential Abundance Testing and Ratio Compression Audit}
# Empirical Bayes linear modeling via limma
design <- model.matrix(~ Condition, data = as.data.frame(colData(qf)))
fit    <- lmFit(prot_mat, design)
fit    <- eBayes(fit)
res    <- topTable(fit, coef = "ConditionTreatment", number = Inf)

mean_obs_log2fc <- mean(res$logFC)
mean_obs_fc     <- 2^mean_obs_log2fc

cat(sprintf("Theoretical Biological Fold Change: 4.00 (log2FC = 2.00)\n"))
cat(sprintf("Empirically Observed Fold Change:   %.2f (log2FC = %.2f)\n", 
            mean_obs_fc, mean_obs_log2fc))
cat(sprintf("Ratio Compression Magnitude:       %.1f%% signal dampening!\n", 
            (1 - mean_obs_fc / 4.0) * 100))
\end{rchunk}

\begin{routput}
Theoretical Biological Fold Change: 4.00 (log2FC = 2.00)
Empirically Observed Fold Change:   3.22 (log2FC = 1.69)
Ratio Compression Magnitude:       19.5\% signal dampening!
\end{routput}

\begin{biologicalinsight}
\textbf{The Clinical Risk of Uncorrected Ratio Compression:}\\
In biomarker discovery experiments, a biological threshold of 2-fold change ($|\log_2\text{FC}| \ge 1.0$) is commonly mandated to select candidates for targeted clinical assay development. Under 50\% co-isolation interference ($f = 0.5$), a true 2.5-fold biological upregulation is observed as:
$$R^{\text{obs}} = 2.5(1 - 0.5) + 0.5 = 1.75 \quad (\log_2\text{FC} = 0.81)$$
Consequently, potent clinical biomarkers are falsely discarded as sub-threshold noise if analysts apply arbitrary fold-change filters to compressed isobaric data without considering precursor ion purity \cite{ting2011sps}.
\end{biologicalinsight}

% =============================================================================
\section{The Five Deadly Sins of Isobaric Labeling Analysis}
% =============================================================================

\begin{pitfallbox}
\textbf{Deadly Sin 1: Interpreting Compressed Fold Changes at Face Value.}\\
\textit{Mechanism:} Mistaking observed TMT fold changes for absolute physiological effect sizes. In isobaric MS2 workflows, co-isolated peptide fragments compress ratios toward 1.0. A measured 1.4-fold change in TMT may reflect a true 3.0-fold biological alteration.\\
\textit{Remedy:} Report moderated $p$-values and FDR alongside fold changes; apply precursor ion purity filtering ($>50\%$), or validate critical candidates using targeted PRM or SPS-MS3 \cite{ting2011sps}.
\end{pitfallbox}

\begin{pitfallbox}
\textbf{Deadly Sin 2: Combining Multiplexes Without a Pooled Internal Reference Channel.}\\
\textit{Mechanism:} Running multiple TMT batches without an identical bridge channel. Cross-multiplex batch effects (differing ionization efficiencies, labeling variances) cannot be decoupled from biological phenotype differences, irreversibly confounding the study.\\
\textit{Remedy:} Always include an identical pooled reference standard (IRS) in channel 1 or 126 across every single multiplex in the cohort \cite{plubell2017extended}.
\end{pitfallbox}

\begin{pitfallbox}
\textbf{Deadly Sin 3: Failing to Apply Isotopic Impurity Correction.}\\
\textit{Mechanism:} Commercially manufactured TMT reagents contain small fractions ($0.5\%$--$2.0\%$) of isotopic impurities ($^{13}\text{C}$ vs. $^{12}\text{C}$ or $^{15}\text{N}$ vs. $^{14}\text{N}$). These spill over into adjacent $\pm 1$ and $\pm 2$ Da channels, introducing systematic bias between adjacent conditions.\\
\textit{Remedy:} Invert the manufacturer lot-specific isotopic spillover matrix $\mathbf{M}_{\text{purity}}$ to deconvolve the raw channel intensities: $\mathbf{I}_{\text{true}} = \mathbf{M}_{\text{purity}}^{-1} \mathbf{I}_{\text{obs}}$.
\end{pitfallbox}

\begin{pitfallbox}
\textbf{Deadly Sin 4: Unweighted Summation of Low-Intensity Reporter Ions.}\\
\textit{Mechanism:} Summing low-intensity reporter ions near the detector noise floor with high-intensity PSMs. Low-abundance reporter ions suffer from Poisson shot-noise and ratio distortion, skewing the aggregate protein ratio.\\
\textit{Remedy:} Use robust weighted M-estimation (\texttt{robustSummary}) or minimum reporter ion intensity thresholds ($>10^3$ counts) \cite{gatto2023qfeatures}.
\end{pitfallbox}

\begin{pitfallbox}
\textbf{Deadly Sin 5: Applying Standard RNA-Seq Count Normalization (e.g., DESeq2).}\\
\textit{Mechanism:} Treating isobaric reporter ion intensities as discrete integer Poisson counts. Reporter intensities represent continuous analog current measurements from an electron multiplier, exhibiting log-normal, heteroscedastic continuous variance structures.\\
\textit{Remedy:} Model isobaric data via linear models with empirical Bayes moderation (\texttt{limma}) or specialized linear mixed-effects packages (\texttt{MSstatsTMT}) \cite{huang2020msstatstmt}.
\end{pitfallbox}

% =============================================================================
\section{Chapter Exercises}
% =============================================================================

\noindent\textbf{Exercise 15.1 (Mathematical Derivation of the De-Compression Bound):}\\
A peptide with true biological ratio $R^{\text{true}} = 8.0$ is analyzed by MS2-TMT. The precursor isolation window contains $60\%$ background interference ($f = 0.60$).
\begin{enumerate}
    \item Calculate the observed fold change $R^{\text{obs}}$ and observed $\log_2\text{FC}$.
    \item If the laboratory instrument measures $R^{\text{obs}} = 3.50$ with empirical standard deviation $\sigma_{\text{obs}} = 0.20$, compute the de-compressed true ratio $\hat{R}^{\text{true}}$ and its standard deviation $\sigma_{\text{true}}$ using the delta method.
\end{enumerate}

\vspace{0.8em}
\noindent\textbf{Exercise 15.2 (Internal Reference Scaling Proof):}\\
Prove that for a three-multiplex study ($M = 3$), multiplying each protein's channel intensity by the IRS scaling factor $S_{gm} = \bar{R}_g / R_{gm}$ guarantees that the scaled reference abundance across all three multiplexes is identical to the global geometric mean $\bar{R}_g$. Show that this transformation is invariant to arbitrary multiplicative gain shifts applied to individual multiplexes.

\vspace{0.8em}
\noindent\textbf{Exercise 15.3 (Isotopic Spillover Matrix Inversion):}\\
Suppose a 3-channel isobaric tag exhibits $1\%$ spillover into the $+1$ adjacent channel and $0.5\%$ spillover into the $-1$ adjacent channel.
\begin{enumerate}
    \item Write the $3 \times 3$ isotopic contamination matrix $\mathbf{M}$.
    \item Derive the inverse matrix $\mathbf{M}^{-1}$ to first-order Taylor approximation.
\end{enumerate}

\vspace{0.8em}
\noindent\textbf{Exercise 15.4 (Computational Audit of Channel CV):}\\
Write an R script that ingests a 16-channel TMT matrix, computes the channel-level coefficient of variation (CV) for invariant housekeeping proteins, and flags channels exceeding a $10\%$ technical noise threshold.

% =============================================================================
\section{Comprehensive Worked-Out Solutions}
% =============================================================================

\subsection*{Solution to Exercise 15.1: Mathematical De-Compression and Delta Method}
\textbf{1. Observed Ratio Under 60\% Interference:}\\
Given $R^{\text{true}} = 8.0$ and $f = 0.60$, from Eq.~(\ref{eq:ratio_compression_linear}):
\begin{align*}
R^{\text{obs}} &= R^{\text{true}}(1 - f) + f = 8.0(1 - 0.60) + 0.60 = 8.0(0.40) + 0.60 = 3.20 + 0.60 = 3.80 \\
\log_2 R^{\text{obs}} &= \log_2(3.80) \approx 1.926
\end{align*}
\textbf{Conclusion:} A dramatic 8-fold biological difference is severely compressed to 3.8-fold ($52.5\%$ ratio reduction).

\textbf{2. De-Compressed Ratio and Variance Inflation:}\\
Given $R^{\text{obs}} = 3.50$ and $\sigma_{\text{obs}} = 0.20$:
\begin{equation*}
\hat{R}^{\text{true}} = \frac{R^{\text{obs}} - f}{1 - f} = \frac{3.50 - 0.60}{1 - 0.60} = \frac{2.90}{0.40} = 7.25
\end{equation*}
Using Eq.~(\ref{eq:variance_inflation}) for the standard deviation:
\begin{equation*}
\sigma_{\text{true}} = \frac{\sigma_{\text{obs}}}{1 - f} = \frac{0.20}{0.40} = 0.50
\end{equation*}
\textbf{Interpretation:} The de-compressed biological fold change is $7.25 \pm 0.50$. While the point estimate is corrected, the uncertainty is magnified by $2.5\times$, illustrating why hardware-level interference elimination (SPS-MS3) is superior to post-hoc mathematical deconvolution.

\subsection*{Solution to Exercise 15.2: Invariance of Internal Reference Scaling}
Let $R_{g1}, R_{g2}, R_{g3}$ be the reference channel intensities in multiplexes $m = 1, 2, 3$. The geometric mean is:
$$\bar{R}_g = \left( R_{g1} R_{g2} R_{g3} \right)^{1/3}$$
The scaled reference value for multiplex $m$ is:
$$\tilde{R}_{gm} = R_{gm} \cdot S_{gm} = R_{gm} \cdot \frac{\bar{R}_g}{R_{gm}} = \bar{R}_g$$
Because $\tilde{R}_{gm} = \bar{R}_g$ identically for all $m \in \{1, 2, 3\}$, the reference signal is perfectly equalized across all multiplexes.

Now suppose multiplex 2 experiences an arbitrary technical gain shift $k > 0$, transforming all intensities in that run such that $y_{g2i}' = k \cdot y_{g2i}$ and $R_{g2}' = k \cdot R_{g2}$. The new geometric mean is:
$$\bar{R}_g' = \left( R_{g1} \cdot (k R_{g2}) \cdot R_{g3} \right)^{1/3} = k^{1/3} \bar{R}_g$$
The new scaling factor for multiplex 2 is:
$$S_{g2}' = \frac{\bar{R}_g'}{R_{g2}'} = \frac{k^{1/3} \bar{R}_g}{k R_{g2}} = k^{-2/3} S_{g2}$$
The scaled intensity in multiplex 2 becomes:
$$\tilde{y}_{g2i}' = y_{g2i}' \cdot S_{g2}' = (k y_{g2i}) \cdot (k^{-2/3} S_{g2}) = k^{1/3} \left( y_{g2i} S_{g2} \right) = k^{1/3} \tilde{y}_{g2i}$$
Simultaneously, for multiplex 1 (unscaled by $k$):
$$S_{g1}' = \frac{\bar{R}_g'}{R_{g1}} = \frac{k^{1/3} \bar{R}_g}{R_{g1}} = k^{1/3} S_{g1} \implies \tilde{y}_{g1i}' = y_{g1i} \cdot (k^{1/3} S_{g1}) = k^{1/3} \tilde{y}_{g1i}$$
Taking the ratio between any sample in multiplex 2 and multiplex 1:
$$\frac{\tilde{y}_{g2i}'}{\tilde{y}_{g1j}'} = \frac{k^{1/3} \tilde{y}_{g2i}}{k^{1/3} \tilde{y}_{g1j}} = \frac{\tilde{y}_{g2i}}{\tilde{y}_{g1j}}$$
The arbitrary multiplicative gain factor $k$ cancels exactly, proving complete invariance to batch sensitivity shifts \cite{plubell2017extended}.

\subsection*{Solution to Exercise 15.3: Isotopic Spillover Matrix Inversion}
\textbf{1. Contamination Matrix:}\\
The spillover matrix $\mathbf{M}$ where column $j$ indicates where channel $j$ signal distributes:
\begin{equation*}
\mathbf{M} = \begin{pmatrix}
0.985 & 0.005 & 0.000 \\
0.010 & 0.985 & 0.005 \\
0.000 & 0.010 & 0.990
\end{pmatrix}
\end{equation*}
\textbf{2. First-Order Taylor Inversion:}\\
Writing $\mathbf{M} \approx \mathbf{I} + \mathbf{E}$, where $\mathbf{E}$ represents the off-diagonal error matrix:
$$\mathbf{M}^{-1} \approx \mathbf{I} - \mathbf{E} = \begin{pmatrix}
1.000 & -0.005 & 0.000 \\
-0.010 & 1.000 & -0.005 \\
0.000 & -0.010 & 1.000
\end{pmatrix}$$
Multiplying $\mathbf{M}^{-1} \mathbf{I}^{\text{obs}}$ effectively subtracts $0.5\%$ of channel 2 from channel 1, and $1.0\%$ of channel 1 from channel 2, recovering uncorrupted precursor counts.

\subsection*{Solution to Exercise 15.4: Computational Channel QC in R}
\begin{rchunk}{Automated TMT Channel Coefficient of Variation Audit}
audit_tmt_channels <- function(tmt_matrix, cv_threshold = 10.0) {
  # tmt_matrix: proteins (rows) x channels (columns)
  cv_results <- data.frame(
    Channel = colnames(tmt_matrix),
    Mean_Intensity = colMeans(tmt_matrix, na.rm = TRUE),
    SD_Intensity   = apply(tmt_matrix, 2, sd, na.rm = TRUE)
  ) %>%
    mutate(
      CV_Percent = (SD_Intensity / Mean_Intensity) * 100,
      Status     = ifelse(CV_Percent > cv_threshold, "FLAGGED (High CV)", "PASS")
    )
  
  return(cv_results)
}

# Run on simulated 6-channel data
test_data <- matrix(rnorm(600, mean = 20000, sd = 1500), nrow = 100, ncol = 6)
# Introduce technical failure in channel 4
test_data[, 4] <- test_data[, 4] * rnorm(100, mean = 1, sd = 0.35)
colnames(test_data) <- paste0("TMT_", 1:6)

qc_summary <- audit_tmt_channels(test_data)
print(qc_summary[, c("Channel", "CV_Percent", "Status")])
\end{rchunk}

\begin{routput}
  Channel CV\_Percent            Status
1   TMT\_1   7.512492              PASS
2   TMT\_2   7.481203              PASS
3   TMT\_3   7.530914              PASS
4   TMT\_4  35.210941 FLAGGED (High CV)
5   TMT\_5   7.491823              PASS
6   TMT\_6   7.520112              PASS
\end{routput}

% =============================================================================
\section*{Canonical References}
\addcontentsline{toc}{section}{Canonical References}
% =============================================================================

\begin{thebibliography}{99}

\bibitem{boland2023health}
Boland, M.~R. (2023). \textit{Health Data Analytics with R: A Hands-On Guide}. Springer Nature, Cham, Switzerland. \doi{10.1007/978-3-031-18099-6}.

\bibitem{gatto2023qfeatures}
Gatto, L., Vanderaa, C., \& Breckels, L.~M. (2023). A foundation for reproducible single-cell and bulk proteomics data analysis in Bioconductor. \textit{Nature Communications}, 14, 5698. \doi{10.1038/s41467-023-41556-2}.

\bibitem{huang2020msstatstmt}
Huang, T., Choi, M., Tzouros, M., Golling, S., Pandya, N.~J., Banfai, B., Dunkley, T., \& Vitek, O. (2020). MSstatsTMT: statistical detection of differentially abundant proteins in experiments with isobaric labeling and multiple mixtures. \textit{Molecular \& Cellular Proteomics}, 19(10), 1706--1723. \doi{10.1074/mcp.RA120.002105}.

\bibitem{mcalister2014multinotch}
McAlister, G.~C., Nusinow, D.~P., Jedrychowski, M.~P., W{\"u}hr, M., Huttlin, E.~L., Erickson, B.~K., Rad, R., Haas, W., \& Gygi, S.~P. (2014). MultiNotch MS3 enables accurate, sensitive, and multiplexed detection of differential expression without human intervention. \textit{Analytical Chemistry}, 86(14), 7150--7158. \doi{10.1021/ac502040v}.

\bibitem{plubell2017extended}
Plubell, D.~L., Wilmarth, P.~A., Zhao, Y., Fenton, A.~M., Minnier, J., Gough, P.~J., \dots \& Pamir, N. (2017). Extended multiplexing of tandem mass tags (TMT) labeling reveals age-related protein alterations in mouse aorta. \textit{Molecular \& Cellular Proteomics}, 16(5), 873--890. \doi{10.1074/mcp.M116.065524}.

\bibitem{ritchie2015limma}
Ritchie, M.~E., Phipson, B., Wu, D., Hu, Y., Law, C.~W., Shi, W., \& Smyth, G.~K. (2015). limma powers differential expression analyses for RNA-sequencing and microarray studies. \textit{Nucleic Acids Research}, 43(7), e47. \doi{10.1093/nar/gkv007}.

\bibitem{ross2004multiplexed}
Ross, P.~L., Huang, Y.~N., Marchese, J.~N., Williamson, B., Parker, K., Pappin, D., \dots \& Pappin, D.~J. (2004). Multiplexed protein quantitation in Saccharomyces cerevisiae using amine-reactive isobaric tagging reagents. \textit{Molecular \& Cellular Proteomics}, 3(12), 1154--1169. \doi{10.1074/mcp.M400129-MCP200}.

\bibitem{smyth2004linear}
Smyth, G.~K. (2004). Linear models and empirical Bayes methods for assessing differential expression in microarray experiments. \textit{Statistical Applications in Genetics and Molecular Biology}, 3(1), Article 3. \doi{10.2202/1544-6115.1027}.

\bibitem{thompson2003tandem}
Thompson, A., Sch{\"a}fer, J., Kuhn, K., Kienle, S., Schwarz, J., Schmidt, G., Neumann, T., Johnstone, R., Mohammed, A.~K., \& Hamon, C. (2003). Tandem mass tags: a novel quantification strategy for comparative analysis of complex protein mixtures by MS/MS. \textit{Analytical Chemistry}, 75(8), 1895--1904. \doi{10.1021/ac0262560}.

\bibitem{ting2011sps}
Ting, L., Rad, R., Gygi, S.~P., \& Haas, W. (2011). MS3 eliminates ratio distortion in isobaric multiplexed quantitative proteomics. \textit{Nature Methods}, 8(11), 937--940. \doi{10.1038/nmeth.1714}.

\end{thebibliography}

\end{document}
